\section*{第 3 章}
\addcontentsline{toc}{section}{第 3 章}

\exercise{习题 3.1}
\textbf{题目：}
设
\[
Y(t)=X(t)\cos(2\pi f_c t+\theta),
\]
其中 $X(t)$ 是零均值平稳过程，其自相关函数为 $R_X(\tau)$，功率谱密度为 $P_X(f)$，$\theta$ 是与 $X(t)$ 独立的随机变量。求 $Y(t)$ 的均值、平均自相关函数与功率谱密度。

\par\textbf{解：}
由于 $X(t)$ 的均值为零，且 $\theta$ 与 $X(t)$ 独立，
\[
\mathbb{E}[Y(t)]=\mathbb{E}[X(t)]\mathbb{E}[\cos(2\pi f_c t+\theta)]=0.
\]
平均自相关函数为
\begin{align*}
\overline{R_Y(\tau)}
&=\mathbb{E}[Y(t+\tau)Y(t)] \\
&=\mathbb{E}[X(t+\tau)X(t)]\,
\mathbb{E}\!\left[\cos(2\pi f_c t+2\pi f_c\tau+\theta)\cos(2\pi f_c t+\theta)\right] \\
&=\frac12 R_X(\tau)\cos(2\pi f_c\tau).
\end{align*}
由频移性质，功率谱密度为
\[
P_Y(f)=\frac14 P_X(f-f_c)+\frac14 P_X(f+f_c).
\]

\medskip
\exercise{习题 3.2}
\textbf{题目：}
功率谱密度为 $N_0/2$ 的加性白高斯噪声 $n_w(t)$ 通过滤波器后的输出为 $n(t)$。已知
\[
H(f)=
\begin{cases}
\sqrt{\dfrac{T_s}{2}\bigl(1+\cos\pi fT_s\bigr)},& |f|\le \dfrac{1}{T_s},\\[6pt]
0,& |f|>\dfrac{1}{T_s}.
\end{cases}
\]
求 $n(t)$ 的功率谱密度及平均功率。

\par\textbf{解：}
输出功率谱密度为
\[
P_n(f)=\frac{N_0}{2}|H(f)|^2=
\begin{cases}
\dfrac{N_0T_s}{4}\bigl(1+\cos\pi fT_s\bigr),& |f|\le \dfrac{1}{T_s},\\[6pt]
0,& |f|>\dfrac{1}{T_s}.
\end{cases}
\]
平均功率为
\[
P_n=\int_{-\infty}^{\infty}P_n(f)\,df
=\int_{-1/T_s}^{1/T_s}\frac{N_0T_s}{4}\bigl(1+\cos\pi fT_s\bigr)\,df
=\frac{N_0}{2}.
\]

\medskip
\exercise{习题 3.3}
\textbf{题目：}
功率谱密度为 $N_0/2$ 的加性白高斯噪声 $n_w(t)$ 先通过一个中心频率为 $f_c$、带宽为 $B$ 的理想带通滤波器，然后经过一个微分器，输出为 $y(t)$。求 $y(t)$ 及其同相分量 $y_c(t)$、正交分量 $y_s(t)$ 的功率谱密度与功率。

\par\textbf{解：}
带通滤波器输出 $n(t)$ 的功率谱密度为
\[
P_n(f)=\frac{N_0}{2}|H(f)|^2=
\begin{cases}
\dfrac{N_0}{2},& |f\pm f_c|\le \dfrac{B}{2},\\[6pt]
0,& \text{其他}.
\end{cases}
\]
微分器等效于传递函数 $j2\pi f$，因此
\[
P_y(f)=P_n(f)|j2\pi f|^2=
\begin{cases}
2N_0\pi^2f^2,& |f\pm f_c|\le \dfrac{B}{2},\\[6pt]
0,& \text{其他}.
\end{cases}
\]
同相分量与正交分量有相同的功率谱密度：
\[
P_c(f)=P_s(f)=
\begin{cases}
P_y(f+f_c)+P_y(f-f_c),& |f|<f_c,\\
0,& \text{其他},
\end{cases}
\]
于是
\[
P_c(f)=P_s(f)=
\begin{cases}
4N_0\pi^2(f^2+f_c^2),& |f|\le \dfrac{B}{2},\\[6pt]
0,& \text{其他}.
\end{cases}
\]
相应功率满足
\[
P_y=\int_{-\infty}^{\infty}P_y(f)\,df,\qquad
P_c=P_s=\int_{-\infty}^{\infty}P_c(f)\,df.
\]
又因带通信号同相、正交分量等功率分配，
\[
P_y=P_c+P_s=2P_c=2P_s.
\]

\medskip
\exercise{习题 3.4}
\textbf{题目：}
设
\[
n(t)=n_c(t)\cos(2\pi f_ct)-n_s(t)\sin(2\pi f_ct)
\]
是窄带平稳高斯噪声，其方差为 $\sigma^2$。$n(t)$ 通过“乘以 $2\cos(2\pi f_ct)$ 再经理想低通滤波器”的系统后输出为 $y(t)$。假设 $n_c(t)$ 与 $n_s(t)$ 的带宽等于理想低通滤波器的带宽。求 $y(t)$ 的一维概率密度函数 $p(y)$。

\par\textbf{解：}
该系统等效提取出噪声的同相分量，因此
\[
y(t)=n_c(t).
\]
而 $n_c(t)$ 是零均值高斯过程，其功率为 $\sigma^2$，故
\[
p(y)=\frac{1}{\sqrt{2\pi\sigma^2}}\exp\!\left(-\frac{y^2}{2\sigma^2}\right).
\]

\medskip
\exercise{习题 3.5}
\textbf{题目：}
功率谱密度为 $N_0/2$ 的加性白高斯噪声 $n_w(t)$ 通过截止频率为 $f_H$ 的理想低通滤波器得到 $n(t)$。以速率 $2f_H$ 对 $n(t)$ 进行均匀采样得到序列 $\{n_k\}$，其中
\[
n_k=n\!\left(\frac{k}{2f_H}\right),\qquad k=0,\pm1,\pm2,\dots
\]
求序列 $\{n_k\}$ 的自相关函数 $R_n(m)=\mathbb{E}[n_{k+m}n_k]$，并求 $n_1,n_2,\dots,n_N$ 的联合概率密度函数。

\par\textbf{解：}
滤波后噪声的功率谱密度为
\[
P_n(f)=\frac{N_0}{2}\rect\!\left(\frac{f}{2f_H}\right),
\]
故自相关函数为
\[
\mathbb{E}[n(t+\tau)n(t)]=N_0f_H\,\sinc(2f_H\tau).
\]
于是
\begin{align*}
R_n(m)
&=\mathbb{E}[n_{k+m}n_k]
=\mathbb{E}\!\left[n\!\left(\frac{k+m}{2f_H}\right)n\!\left(\frac{k}{2f_H}\right)\right] \\
&=N_0f_H\,\sinc(m)
=
\begin{cases}
N_0f_H,& m=0,\\
0,& m=\pm1,\pm2,\dots
\end{cases}
\end{align*}
因此各采样值都是零均值高斯随机变量，且方差为 $N_0f_H$；不同采样值彼此不相关，从而相互独立。故
\[
p(n_1,n_2,\dots,n_N)
=\prod_{k=1}^N\frac{1}{\sqrt{2\pi N_0f_H}}
\exp\!\left(-\frac{n_k^2}{2N_0f_H}\right)
\]
\[
=(2\pi N_0f_H)^{-N/2}\exp\!\left(-\frac{1}{2N_0f_H}\sum_{k=1}^N n_k^2\right).
\]

\medskip
\exercise{习题 3.6}
\textbf{题目：}
题中确定信号 $b(t)$ 为幅度仅取 $\pm A$ 的矩形波形：
\[
b(t)=
\begin{cases}
A,& 0\le t<3T_c,\\
-A,& 3T_c\le t<5T_c,\\
A,& 5T_c\le t<6T_c,\\
-A,& 6T_c\le t<7T_c,\\
0,& \text{其他},
\end{cases}
\]
记总持续时间 $T=7T_c$。该信号叠加功率谱密度为 $N_0/2$ 的白高斯噪声后，经过对 $b(t)$ 匹配的匹配滤波器。求匹配滤波器冲激响应、输出端最大瞬时信噪比，以及最佳采样时刻采样值的概率密度。

\par\textbf{解：}
匹配滤波器冲激响应通式为
\[
h(t)=K\,b(t_0-t).
\]
考虑因果性，取最佳采样时刻 $t_0=T$，并令 $K=1$，则
\[
h(t)=b(T-t).
\]
信号能量为
\[
E_b=\int_{-\infty}^{\infty}b^2(t)\,dt=A^2T.
\]
因此匹配滤波器输出端最大瞬时信噪比为
\[
\gamma_{\max}=\frac{2E_b}{N_0}=\frac{2A^2T}{N_0}.
\]
最佳采样时刻输出噪声功率为
\[
\frac{N_0}{2}E_b=\frac{N_0A^2T}{2},
\]
而采样值中的信号分量为
\[
E_b=A^2T.
\]
故采样值可写成
\[
y=A^2T+z,
\]
其中 $z$ 是零均值高斯随机变量，其方差为 $N_0A^2T/2$。因此
\[
p(y)=\frac{1}{\sqrt{\pi N_0A^2T}}
\exp\!\left[-\frac{(y-A^2T)^2}{N_0A^2T}\right].
\]

\medskip
\exercise{习题 3.7}
\textbf{题目：}
设 $X_1,X_2$ 分别服从瑞利分布和莱斯分布，且相互独立。其概率密度函数为
\[
p_{X_1}(x_1)=\frac{x_1}{\sigma^2}e^{-x_1^2/(2\sigma^2)},\qquad x_1\ge 0,
\]
\[
p_{X_2}(x_2)=\frac{x_2}{\sigma^2}e^{-(x_2^2+A^2)/(2\sigma^2)}I_0\!\left(\frac{Ax_2}{\sigma^2}\right),\qquad x_2\ge 0.
\]
求概率 $P(X_1>X_2)$。

\par\textbf{解：}
由独立性，
\[
P(X_1>X_2)=\int_0^\infty p_{X_2}(x_2)\left(\int_{x_2}^{\infty}p_{X_1}(x_1)\,dx_1\right)dx_2.
\]
其中
\[
\int_{x_2}^{\infty}p_{X_1}(x_1)\,dx_1
=\int_{x_2}^{\infty}\frac{x_1}{\sigma^2}e^{-x_1^2/(2\sigma^2)}\,dx_1
=e^{-x_2^2/(2\sigma^2)}.
\]
代回得
\[
P(X_1>X_2)=\int_0^\infty \frac{x_2}{\sigma^2}
e^{-(2x_2^2+A^2)/(2\sigma^2)}I_0\!\left(\frac{Ax_2}{\sigma^2}\right)\,dx_2.
\]
令
\[
x=\sqrt{2}\,x_2,\qquad A'=\frac{A}{\sqrt2},
\]
则上式化为
\[
P(X_1>X_2)=\frac12 e^{-A^2/(2\sigma^2)}
\int_0^\infty \frac{x}{\sigma^2}
e^{-(x^2+A'^2)/(2\sigma^2)}I_0\!\left(\frac{A'x}{\sigma^2}\right)\,dx.
\]
最后一个积分是莱斯概率密度函数的积分，等于 $1$，故
\[
P(X_1>X_2)=\frac12 e^{-A^2/(2\sigma^2)}.
\]

\medskip
\exercise{习题 3.8}
\textbf{题目：}
加性白高斯噪声 $n_w(t)$ 通过滤波器 $H_1(f)$ 和 $H_2(f)$ 后的输出分别为 $n_1(t)$ 和 $n_2(t)$。问何种 $H_1(f)$ 和 $H_2(f)$ 可保证 $n_1(t),n_2(t)$ 统计独立？

\par\textbf{解：}
设对应冲激响应分别为 $h_1(t),h_2(t)$。则
\[
n_1(t)=\int_{-\infty}^{\infty}h_1(u)n_w(t-u)\,du,
\qquad
n_2(t)=\int_{-\infty}^{\infty}h_2(v)n_w(t-v)\,dv.
\]
统计独立要求它们的互相关函数恒为零：
\begin{align*}
R_{12}(\tau)
&=\mathbb{E}[n_1(t+\tau)n_2(t)] \\
&=\frac{N_0}{2}\int_{-\infty}^{\infty}h_1(u)h_2(u-\tau)\,du.
\end{align*}
最后一个积分就是 $h_1(t)$ 与 $h_2(t)$ 的互相关函数。若它恒为零，则互能量谱密度为零，即
\[
H_1(f)H_2^*(f)=0,\qquad -\infty<f<\infty.
\]
因此，当且仅当 $H_1(f)$ 与 $H_2(f)$ 在频域中互不重叠时，$n_1(t)$ 与 $n_2(t)$ 统计独立。

\medskip
\exercise{习题 3.9}
\textbf{题目：}
设 $X(t)$ 是零均值平稳高斯过程，其自相关函数为
\[
R_X(\tau)=\sinc^2\!\left(\frac{\tau}{T}\right).
\]
令
\[
Y(t)=X(t-T).
\]
求 $Y(t)$ 的自相关函数、$Y(t)$ 与 $X(t)$ 的互相关函数、$Y(0)$ 与 $X(0)$ 的联合概率密度函数，以及 $Y(T)$ 与 $X(0)$ 的联合概率密度函数。

\par\textbf{解：}
$Y(t)$ 是 $X(t)$ 的延迟，因此仍为零均值平稳高斯过程，且延迟不改变自相关函数：
\[
R_Y(\tau)=R_X(\tau)=\sinc^2\!\left(\frac{\tau}{T}\right).
\]
互相关函数为
\[
R_{YX}(\tau)=\mathbb{E}[Y(t+\tau)X(t)]
=\mathbb{E}[X(t+\tau-T)X(t)]
=R_X(\tau-T)
=\sinc^2\!\left(\frac{\tau}{T}-1\right).
\]
当 $\tau=0$ 时，
\[
R_{YX}(0)=\sinc^2(-1)=0,
\]
故 $Y(0)$ 与 $X(0)$ 不相关。由于它们是联合高斯，故相互独立；又各自方差均为 $R_X(0)=1$，于是
\[
p_{Y(0),X(0)}(y,x)=\frac{1}{2\pi}\exp\!\left(-\frac{y^2+x^2}{2}\right).
\]
另一方面，
\[
Y(T)=X(T-T)=X(0),
\]
即 $Y(T)$ 与 $X(0)$ 是同一个随机变量，因此
\[
p_{Y(T),X(0)}(y,x)=\frac{1}{\sqrt{2\pi}}e^{-x^2/2}\delta(y-x).
\]

\medskip
\exercise{习题 3.10}
\textbf{题目：}
确定能量信号 $s(t)$ 的自相关函数为
\[
R_s(\tau)=4\sinc^2\!\left(\frac{\tau}{T}\right).
\]
$s(t)$ 叠加了功率谱密度为 $N_0/2$ 的白高斯噪声 $n_w(t)$ 后，通过一个冲激响应为
\[
g(t)=s(-t)
\]
的滤波器。求滤波器输出端任意时刻 $t$ 的瞬时信噪比。

\par\textbf{解：}
滤波器能量等于信号能量，即
\[
E_g=E_s=R_s(0)=4.
\]
输出噪声 $z(t)$ 是零均值平稳过程，其方差为
\[
\frac{N_0}{2}E_g=2N_0.
\]
信号分量通过滤波器后的输出为
\[
y(t)=\int_{-\infty}^{\infty}g(\tau)s(t-\tau)\,d\tau
=\int_{-\infty}^{\infty}s(-\tau)s(t-\tau)\,d\tau
=R_s(t)
=4\sinc^2\!\left(\frac{t}{T}\right).
\]
因此输出端任意时刻的瞬时信噪比为
\[
\gamma(t)=\frac{y^2(t)}{\mathbb{E}[z^2(t)]}
=\frac{8}{N_0}\sinc^4\!\left(\frac{t}{T}\right).
\]
